\documentclass[aspectratio=169,11pt]{beamer}

\usepackage[utf8]{inputenc}
\usepackage[T1]{fontenc}
\usepackage[brazil]{babel}
\usepackage{amsmath,amssymb,amsfonts}
\usepackage{booktabs,tabularx}
\usepackage{graphicx}
\usepackage{xcolor}
\usepackage{hyperref}
\usepackage{listings}
\usepackage{tikz}
\usetikzlibrary{positioning}

\definecolor{ufopred}{RGB}{140,29,36}
\definecolor{deepnavy}{RGB}{20,35,60}
\definecolor{lightgray}{RGB}{245,247,250}
\definecolor{darkslate}{RGB}{30,40,50}
\definecolor{softblue}{RGB}{230,238,247}

\usetheme{Madrid}
\usecolortheme{default}
\setbeamertemplate{navigation symbols}{}
\setbeamercolor{structure}{fg=ufopred}
\setbeamercolor{frametitle}{fg=white,bg=ufopred}
\setbeamercolor{title}{fg=white,bg=ufopred}
\setbeamercolor{block title}{fg=white,bg=deepnavy}
\setbeamercolor{block body}{fg=black,bg=lightgray}
\setbeamercolor{alerted text}{fg=ufopred}
\setbeamertemplate{footline}{%
  \leavevmode\hbox{%
    \begin{beamercolorbox}[wd=.72\paperwidth,ht=2.6ex,dp=1.1ex,leftskip=1.5ex]{author in head/foot}
      \textcolor{darkslate}{\insertshorttitle}
    \end{beamercolorbox}%
    \begin{beamercolorbox}[wd=.28\paperwidth,ht=2.6ex,dp=1.1ex,rightskip=1.5ex plus1fil]{date in head/foot}
      \hfill \textcolor{darkslate}{\insertframenumber/\inserttotalframenumber}
    \end{beamercolorbox}%
  }%
}
\setbeamertemplate{caption}[numbered]
\graphicspath{{figures/}}
\newcommand{\status}[1]{\textcolor{ufopred}{\textbf{#1}}}

% Figura opcional: se o arquivo não existir, mostra um espaço reservado
% uso: \fig{arquivo-sem-extensao}{largura}{altura}
\newcommand{\fig}[3]{%
  \IfFileExists{figures/#1.png}%
    {\includegraphics[width=#2,height=#3,keepaspectratio]{#1.png}}%
    {\fbox{\parbox[c][#3][c]{#2}{\centering\footnotesize\textcolor{ufopred}{[inserir figura: \detokenize{#1}]}}}}}

\title[Visualização CCO -- Computação Gráfica]{Visualização Geométrica e Rasterização de Árvores Arteriais CCO}
\subtitle{BCC327 -- Computação Gráfica -- TP1}
\author{Cauã Lemos Câmara \and César Augusto Tiago Totô \and Roberto Antonino Menegassi Filho}
\institute[UFOP]{Departamento de Computação (DECOM)\\Universidade Federal de Ouro Preto (UFOP)}
\date{11 de novembro de 2026}

\begin{document}

\begin{frame}
  \titlepage
\end{frame}

\begin{frame}{Roteiro}
  \tableofcontents
\end{frame}

% =========================================================
\section{Introdução e Objetivos}
% =========================================================
\begin{frame}{Contexto e Objetivo}
\begin{block}{Problema}
Como interpretar e visualizar árvores arteriais geradas pelo método CCO e exercitar conceitos de Computação Gráfica?
\end{block}
\begin{itemize}
  \item Modelos CCO fornecidos pelo docente em arquivos \texttt{.vtk}.
  \item O método CCO \alert{não} foi reimplementado.
  \item Foco: leitura, representação, rasterização e recorte.
  \item Distinguir o que é \alert{implementação do grupo} do que é \alert{biblioteca}.
\end{itemize}
\end{frame}

\begin{frame}{Objetivos Específicos}
\begin{enumerate}
  \item Ler a estrutura dos modelos (pontos, conectividade, raio).
  \item Representar a árvore por tubos, com raio em cor e espessura.
  \item Explorar vistas e projeções (ortográfica e perspectiva).
  \item Rasterizar linhas (DDA, Bresenham) e polígonos (Scanline).
  \item Comparar renderização com e sem antialiasing.
  \item Recortar segmentos (Cohen-Sutherland, Liang-Barsky) e polígonos (Sutherland-Hodgman).
\end{enumerate}
\end{frame}

% =========================================================
\section{Fundamentação Teórica}
% =========================================================
\begin{frame}{Modelo CCO e Codificação do Raio}
\begin{columns}[T,onlytextwidth]
\column{0.52\textwidth}
\begin{itemize}
  \item VTK POLYDATA: pontos, células lineares e raio por célula.
  \item Árvore = vértices + arestas com raio.
  \item Tubos gerados pelo filtro \texttt{tube} (8 lados, extremidades fechadas).
\end{itemize}
\column{0.44\textwidth}
\begin{block}{Cor pelo raio}
\[ t=\frac{r-r_{\min}}{r_{\max}-r_{\min}} \]
\[ R=t,\quad G=0{,}5,\quad B=1-t \]
\end{block}
Raio menor $\rightarrow$ azul; raio maior $\rightarrow$ vermelho.
\end{columns}
\end{frame}

\begin{frame}{Câmera, Vistas e Projeção}
\begin{columns}[T,onlytextwidth]
\column{0.48\textwidth}
\begin{block}{Perspectiva}
Objetos distantes aparecem menores.
\end{block}
\column{0.48\textwidth}
\begin{block}{Ortográfica (paralela)}
Paralelas se mantêm; tamanho independe da distância.
\end{block}
\end{columns}
\vspace{0.2cm}
\begin{itemize}
  \item Vistas predefinidas: \texttt{yz}, \texttt{xz} e \texttt{xy}.
  \item Rotação livre e zoom pela interação padrão do PyVista/VTK.
\end{itemize}
\end{frame}

\begin{frame}{Rasterização de Linhas e Polígonos}
\begin{columns}[T,onlytextwidth]
\column{0.48\textwidth}
\begin{block}{Linhas}
\begin{itemize}
  \item \textbf{DDA}: passo em ponto flutuante.
  \item \textbf{Bresenham}: apenas inteiros.
\end{itemize}
\[ p_0 = 2\Delta y - \Delta x \]
\end{block}
\column{0.48\textwidth}
\begin{block}{Scanline Fill}
\begin{itemize}
  \item Tabelas ET e AET.
  \item Regra par--ímpar.
\end{itemize}
\[ x_{k+1}=x_k+\tfrac{1}{m} \]
\end{block}
\end{columns}
\end{frame}

\begin{frame}{Recorte e Antialiasing}
\begin{columns}[T,onlytextwidth]
\column{0.48\textwidth}
\begin{block}{Recorte}
\begin{itemize}
  \item \textbf{Cohen-Sutherland}: códigos de 4 bits.
  \item \textbf{Liang-Barsky}: forma paramétrica.
  \item \textbf{Sutherland-Hodgman}: uma borda por vez.
  \item Cyrus-Beck e Weiler-Atherton: apenas discutidos.
\end{itemize}
\end{block}
\column{0.48\textwidth}
\begin{block}{Antialiasing}
\begin{itemize}
  \item Reduz o serrilhado das bordas.
  \item MSAA, FXAA e SSAA oferecidos pelo VTK.
\end{itemize}
\end{block}
\end{columns}
\end{frame}

% =========================================================
\section{Metodologia e Arquitetura}
% =========================================================
\begin{frame}{Tecnologias e Fluxo}
\begin{columns}[T,onlytextwidth]
\column{0.46\textwidth}
\begin{itemize}
  \item \textbf{Python}: aplicação (\texttt{main.py}).
  \item \textbf{PyVista / VTK}: leitura, tubos, renderização.
  \item \textbf{NumPy}: dados e grade de pixels.
  \item \textbf{Matplotlib}: mapa de cores.
\end{itemize}
\column{0.50\textwidth}
\centering
\begin{tikzpicture}[
  node distance=0.35cm,
  every node/.style={draw,rounded corners,align=center,minimum width=5.6cm,minimum height=0.5cm,font=\footnotesize},
  arrow/.style={->,thick,ufopred}
]
\node (a) {Leitura do VTK};
\node (b) [below=of a] {Pontos, segmentos e raios};
\node (c) [below=of b] {Tubos e atributos visuais};
\node (d) [below=of c] {Câmera, rasterização e recorte};
\draw[arrow] (a)--(b);
\draw[arrow] (b)--(c);
\draw[arrow] (c)--(d);
\end{tikzpicture}
\end{columns}
\end{frame}

\begin{frame}{Módulos da Aplicação}
\small
\begin{tabularx}{\textwidth}{c p{0.34\textwidth} X}
\toprule
\textbf{Mód.} & \textbf{Tema} & \textbf{Situação}\\
\midrule
1 & Leitura e atributos & \status{Implementado}\\
2 & Câmera e vistas & \status{Parcial} (vistas e projeção via API)\\
3 & Rasterização de linhas & \status{Implementado} (DDA e Bresenham)\\
4 & Rasterização de polígonos & \status{Implementado} (Scanline)\\
5 & Antialiasing & \status{Via biblioteca}\\
6 & Recorte de segmentos & \status{Implementado} (CS e LB)\\
7 & Recorte de polígonos & \status{Implementado}\\
\bottomrule
\end{tabularx}
\end{frame}

\begin{frame}{Detalhes de Implementação}
\begin{itemize}
  \item \textbf{Linhas}: pontos da árvore levados à grade inteira; 20 primeiros segmentos rasterizados.
  \item \textbf{Polígonos}: cada segmento vira um retângulo de 4 vértices; buffer de $400\times400$ pixels; ordenação por profundidade (algoritmo do pintor).
  \item \textbf{Recorte}: janela retangular com slider (10\% a 120\% do modelo); contadores de interseções e divisões.
\end{itemize}
\end{frame}

% =========================================================
\section{Resultados}
% =========================================================
\begin{frame}{Modelo e Atributos}
\begin{columns}[T,onlytextwidth]
\column{0.34\textwidth}
\begin{block}{\texttt{tree3D\_Nterm016.vtk}}
\centering
\begin{tabular}{lr}
\toprule
Pontos & 32\\
Segmentos & 31\\
Raio mín. & 0,2261\\
Raio máx. & 0,6870\\
\bottomrule
\end{tabular}
\end{block}
\column{0.62\textwidth}
\centering
\fig{aplicacao_modelo_3d}{\linewidth}{0.66\textheight}
\end{columns}
\end{frame}

\begin{frame}{Comparação com a Referência}
\begin{columns}[T,onlytextwidth]
\column{0.34\textwidth}
\centering
\fig{referencia_professor_3d}{\linewidth}{0.62\textheight}\\
{\footnotesize Referência do docente.}
\column{0.62\textwidth}
\centering
\fig{aplicacao_modelo_3d}{\linewidth}{0.62\textheight}\\
{\footnotesize Aplicação.}
\end{columns}
\vspace{0.1cm}
{\small Estrutura preservada; a aplicação adiciona controles e codificação pelo raio.}
\end{frame}

\begin{frame}{Codificação por Cor e Espessura}
\begin{columns}[T,onlytextwidth]
\column{0.49\textwidth}
\centering
\fig{atributos_cor}{\linewidth}{0.56\textheight}\\
{\footnotesize Cor associada ao raio.}
\column{0.49\textwidth}
\centering
\fig{cor_desativada}{\linewidth}{0.56\textheight}\\
{\footnotesize Cor uniforme.}
\end{columns}
\vspace{0.1cm}
{\small Cor e espessura podem ser ativadas de forma independente.}
\end{frame}

\begin{frame}{Rasterização de Linhas}
\begin{columns}[T,onlytextwidth]
\column{0.40\textwidth}
\begin{itemize}
  \item DDA e Bresenham retornam a lista de pixels $(x,y)$.
  \item Bresenham trata todos os octantes.
  \item Comparação com a reta ideal.
\end{itemize}
\column{0.56\textwidth}
\fig{rasterizacao_linhas}{\linewidth}{0.62\textheight}
\end{columns}
\end{frame}

\begin{frame}{Rasterização de Polígonos: Scanline}
\begin{columns}[T,onlytextwidth]
\column{0.34\textwidth}
\begin{itemize}
  \item Teste em polígono côncavo.
  \item Grade de $30\times26$ pixels.
  \item \textbf{216} de 780 pixels preenchidos (\textbf{27,7\%}).
\end{itemize}
\column{0.62\textwidth}
\centering
\fig{scanline_3d}{\linewidth}{0.66\textheight}\\
{\footnotesize Vista superior com a grade rasterizada.}
\end{columns}
\end{frame}

\begin{frame}{Scanline Aplicado à Árvore 2D}
\centering
\fig{scanline}{0.85\textwidth}{0.58\textheight}\\[0.1cm]
{\footnotesize Esquerda: árvore 2D. Direita: árvore rasterizada (Scanline).}
\end{frame}

\begin{frame}{Vistas Predefinidas}
\begin{columns}[T,onlytextwidth]
\column{0.32\textwidth}
\centering
\fig{vista_frontal}{\linewidth}{0.45\textheight}\\
{\footnotesize Frontal}
\column{0.32\textwidth}
\centering
\fig{vista_lateral}{\linewidth}{0.45\textheight}\\
{\footnotesize Lateral}
\column{0.32\textwidth}
\centering
\fig{vista_superior}{\linewidth}{0.45\textheight}\\
{\footnotesize Superior}
\end{columns}
\vspace{0.15cm}
\begin{alertblock}{Atenção}
Os rótulos dos botões de projeção estão invertidos em relação às chamadas da API. A projeção ativa deve ser interpretada pelo comportamento da câmera, não pelo texto do botão.
\end{alertblock}
\end{frame}

\begin{frame}{Antialiasing}
\begin{columns}[T,onlytextwidth]
\column{0.49\textwidth}
\centering
\fig{aa_sem}{\linewidth}{0.33\textheight}\\
{\footnotesize Sem antialiasing}\\[0.1cm]
\fig{aa_fxaa}{\linewidth}{0.33\textheight}\\
{\footnotesize FXAA}
\column{0.49\textwidth}
\centering
\fig{aa_msaa}{\linewidth}{0.33\textheight}\\
{\footnotesize MSAA}\\[0.1cm]
\fig{aa_ssaa}{\linewidth}{0.33\textheight}\\
{\footnotesize SSAA}
\end{columns}
\end{frame}

\begin{frame}{Antialiasing: Observações}
\begin{itemize}
  \item Modos selecionáveis: nenhum, MSAA, FXAA e SSAA.
  \item Execução interna feita pelo PyVista/VTK, não pelo grupo.
  \item Sem comparação quantitativa: apenas capturas da seleção dos modos.
\end{itemize}
\end{frame}

\begin{frame}{Recorte de Segmentos}
\begin{columns}[T,onlytextwidth]
\column{0.40\textwidth}
\begin{itemize}
  \item Janela em laranja.
  \item Partes mantidas em verde.
  \item Partes descartadas em vermelho.
  \item Painel com segmentos inteiros, cortados e rejeitados.
\end{itemize}
\column{0.56\textwidth}
\fig{recorte}{\linewidth}{0.62\textheight}
\end{columns}
\end{frame}

% =========================================================
\section{Discussão e Conclusão}
% =========================================================
\begin{frame}{Discussão}
\begin{itemize}
  \item \textbf{Próprio do grupo:} leitura VTK, codificação por raio, DDA/Bresenham, Scanline e recorte.
  \item \textbf{Via biblioteca:} tubos, projeção, vistas e antialiasing.
  \item \textbf{Limitações:} sem translação, rotação e escala explícitas; rótulos de projeção invertidos; antialiasing sem comparação quantitativa.
\end{itemize}
\begin{block}{Desafio principal}
Integrar dados, atributos e visualização sem atribuir ao grupo algoritmos executados pela biblioteca.
\end{block}
\end{frame}

\begin{frame}{Conclusão}
\begin{itemize}
  \item Aplicação para visualizar modelos CCO com leitura VTK, tubos e codificação do raio.
  \item Algoritmos clássicos de rasterização e recorte implementados pelo grupo.
  \item Trabalho futuro: controles geométricos explícitos, correção dos rótulos de projeção e comparação controlada de antialiasing.
\end{itemize}
\end{frame}

\begin{frame}[allowframebreaks]{Referências}
\scriptsize
\begin{thebibliography}{9}
\bibitem{schreiner} SCHREINER, W.; BUXBAUM, P. Computer-optimization of vascular trees. \textit{IEEE Trans. Biomed. Eng.}, v. 40, n. 5, p. 482--491, 1993.
\bibitem{foley} FOLEY, J. D. et al. \textit{Computer Graphics: Principles and Practice}. 2. ed. Addison-Wesley, 1996.
\bibitem{pyvista} PYVISTA DEVELOPERS. \textit{PyVista Documentation}. \url{https://docs.pyvista.org/}
\bibitem{vtk} KITWARE. \textit{VTK -- The Visualization Toolkit}. \url{https://vtk.org/}
\bibitem{bresenham} BRESENHAM, J. E. Algorithm for computer control of a digital plotter. \textit{IBM Systems Journal}, v. 4, n. 1, p. 25--30, 1965.
\bibitem{liang} LIANG, Y.-D.; BARSKY, B. A. A new concept and method for line clipping. \textit{ACM TOG}, v. 3, n. 1, p. 1--22, 1984.
\bibitem{cyrus} CYRUS, M.; BECK, J. Generalized two- and three-dimensional clipping. \textit{Computers \& Graphics}, v. 3, n. 1, p. 23--28, 1978.
\bibitem{sutherland} SUTHERLAND, I. E.; HODGMAN, G. W. Reentrant polygon clipping. \textit{Commun. ACM}, v. 17, n. 1, p. 32--42, 1974.
\bibitem{weiler} WEILER, K.; ATHERTON, P. Hidden surface removal using polygon area sorting. \textit{SIGGRAPH '77}, p. 214--222, 1977.
\end{thebibliography}
\end{frame}

\begin{frame}{Obrigado!}
\centering
{\Huge Obrigado!}\\[0.45cm]
{\large Perguntas?}\\[0.4cm]
{\small Cauã Lemos Câmara \quad César Augusto Tiago Totô\\
Roberto Antonino Menegassi Filho}
\end{frame}

\end{document}